\section{相关化合物的合成方法}
\label{sec:neighbor-routes}

本节按相形成过程比较相关化合物的合成方法，而不逐篇罗列文献。图~\ref{fig:route-matrix}\CJKecglue{}概括可比较的实验项目，表~\ref{tab:synthesis-conditions}\CJKecglue{}保留数值和文献出处。以下仅讨论相关研究提供的方法信息，不将任何路线改写为Ba$_5$Y$_{12}$Zn[O(SiO$_4$)]$_8$的候选配方。

\subsection{高温溶液法路线}

Y--Si--O体系的文献首先表明，“获得晶体”不等于“结构归属正确”。以Na$_2$MoO$_4$为助熔剂得到的$\beta$-Y$_2$Si$_2$O$_7$伴有其他硅酸盐晶体，最终物相归属由单晶结构测定确认；此类研究可在营养物与助熔剂配比、坩埚加盖、缓冷和晶体分离方面提供参照，但不能给出目标相的稳定范围\cite{redhammer2003beta}。以$^{89}$Y MAS\nobreakdash-NMR重新考察Y$_2$Si$_2$O$_7$\slash Y$_2$SiO$_5$多晶型的研究表明，局域探针可补充衍射方法以区分多晶型，但不能单独证明目标相的相形成\cite{becerro2004revisiting}。Y--Si--O合成综述和经典的稀土硅酸盐晶体化学分类可统一多晶型与硅酸根聚合度的术语，但不能替代原始实验数据\cite{sun2014recent,felsche1973the}。

在Ba--Zn--Si--O体系中，开放体系高温溶液法研究仅确认了Ba$_3$Zn$_4$Si$_4$O$_{15}$的结构和路线类别，精确批料与热历史仍记为未报道\cite{ababaikeri2024ba3zn4si4o15}。Zn$_2$SiO$_4$的早期熔融溶液生长研究表明，缓冷过程、晶体分离方法和助熔剂组成必须同时给出；含Pb\slash F的条件还涉及毒性、挥发和废物处置审查，不能作为默认选择\cite{giess1982zn2sio4}。来源明确的无坩埚熔体研究和Y$_2$SiO$_5$\slash Y$_2$Si$_2$O$_7$高温结晶研究，可在容器接触、自发成核和结构鉴定方面提供参照，但不能证明目标相具有相同的熔体生长行为\cite{christensen1994investigation,leonyuk1999high}。Leonyuk的另一篇文献仅有题名层面的证据，只能确认其单晶对象和结构精修，不能据此判断合成方法\cite{leonyuk1999crystal}。

\paragraph{对目标实验的启示。}应分别列出投料与助熔剂组成，注明容器接触情况，并完整记录冷却与分离过程；晶体结构、块体衍射和成分分析需配套进行。相关化合物的具体数值、Pb\slash F体系的适用性以及由晶体外形推定目标相归属，均不能直接外推。

\subsection{固相反应路线}

BaMSiO$_4$（M=Co、Zn、Mg）经单晶X射线衍射和粉末中子衍射建立了填充鳞石英型（stuffed tridymite）同构系列，说明同价四面体位取代可保持母体拓扑，但该骨架不能等同于目标相的零维混合构筑\cite{liu1993structures}。单斜Ba$_2$ZnSi$_2$O$_7$的单晶精修和单斜Ba$_2$MgSi$_2$O$_7$的结构报道分别提供了Zn和Mg参照；二者可用于识别竞争相和检查取代风险，但不能据此断言其与目标相同构\cite{kaiser2002crystal,aitasalo2006crystal}。

受控取代系列示范了组成、结构、热历史与表征的配套报道方式。Ba\slash Sr取代物由SCXRD与EDS联合确认；BaZn$_{2-x}$Co$_x$Si$_2$O$_7$以及Ba\slash Sr--Zn--Si\slash Ge系列的HT-XRD结果给出了与组成相关的结构保持区和相边界。这些研究说明，名义取代量必须与实际组成和原位结构结果对应，不能跨体系量化目标相的Y\slash Zn容限\cite{thieme2015ba1,kerstan2013bazn2si2o7,thieme2016negative}。BaZn$_2$Si$_2$O$_7$固溶体综述可梳理热膨胀随组成变化的背景，但不能替代单个原始烧结批次\cite{thieme2022solid}。含Ba$_{1-x}$Sr$_x$Zn$_2$Si$_2$O$_7$的微晶玻璃研究同样只提供功能与晶化背景，不能作为目标相的直接合成依据\cite{thieme2017variable}。多气氛Ba$_2$ZnSi$_2$O$_7$陶瓷的对照研究表明，气氛变化必须结合Zn损失、第二相和表面价态解释，不能仅凭介电性能变化判断结构稳定性\cite{zou2019anti}。

\paragraph{对目标实验的启示。}应并列给出化学计量称量、复磨复烧、气氛和实际组成，并采用PXRD\slash Rietveld、成分分析及必要的HT-XRD表征。相关研究中的同价取代、热膨胀或介电性能结果，不能直接改写为目标相的异价取代机制、相纯窗口或还原稳定性。

\subsection{机械化学预活化与提拉法工艺参照}

机械化学研究仅支持“预活化可能改变后续相形成路径”这一认识。钇氧磷灰石研究给出了空气中256\,rpm、1--3\,h以及玛瑙罐和玛瑙球的机械处理条件；处理后仍需热处理，且多个Y--Si--O相共存。这些数值仅描述该混相体系，不能直接套用为目标相的活化条件\cite{tzvetkov2001effects}。Y$_2$SiO$_5$粉体研究仅明确报道了LiYO$_2$辅助的固--液相烧结以及显微、XRD和SEM--EDS表征，不能因研究对象为粉体而归为机械化学方法\cite{yldrm2009y2sio5}。

提拉法研究提供了熔体、铱坩埚、气氛、SiO$_2$挥发、提拉与旋转以及退火等方面的参照。Y$_2$SiO$_5$的生长、形貌和光谱联合研究说明，晶体生长后仍需进行结构和成分鉴定。稀土正硅酸盐的早期提拉法研究及后续学位论文可用于规范工艺参数的报道，但目前尚无目标相的提拉法生长证据\cite{pang2005study,shoudu1999czochralski,brandle1986czochralski,alizadeh2021spectroscopy}。

\paragraph{对目标实验的启示。}可参照机械化学活化过程中的污染检查，以及提拉法对熔体、坩埚和挥发的控制；同时必须验证块体相组成、实际元素比、夹杂和单晶结构。普通混匀不能称为机械化学活化，Y$_2$SiO$_5$\slash Ln$_2$SiO$_5$的可提拉性也不等同于目标相的可提拉性。所有数值均以表~\ref{tab:synthesis-conditions}\CJKecglue{}所列原始文献数据为限。
